\section{The Fano surface and the geometric glue}\label{sec:geometric-glue}

Let \(F\) be the Fano surface of lines on \(X\), and let
\(a:F\hookrightarrow J\) be its Albanese embedding.  The difference map
\[
 \psi:F\times F\longrightarrow\Theta,
 \qquad (r,s)\longmapsto a(s)-a(r)
\]
is generically finite of degree six.  Blowing up the diagonal gives a
commutative diagram
\[
 \begin{array}{ccc}
 Y=\operatorname{Bl}_{\Delta}(F\times F)&\xrightarrow{q}&M\\
 \downarrow\mu&&\downarrow b\\
 F\times F&\xrightarrow{\psi}&J.
 \end{array}
\]
The exceptional divisor of \(\mu\) is
\(P=\mathbf P(T_F)\), the universal line
\(\{(\ell,x):x\in\ell\}\subset F\times X\).  Write
\(\pi:P\to F\) and \(p:P\to X\) for its projections.  These facts, as
well as \([F]=\theta^{[3]}\), are classical
\cite{ClemensGriffiths,BeauvilleTheta}.

The cylinder map
\[
 \pi_*p^*:H^3(X,\Z)\longrightarrow H^1(F,\Z)
\]
is an integral isomorphism by Clemens--Griffiths
\cite[\S2 and \S11]{ClemensGriffiths}.  Its adjoint for the unimodular
pairings is therefore an isomorphism
\begin{equation}\label{eq:cylinder-adjoint}
 p_*\pi^*:H^3(F,\Z)/\operatorname{tors}
 \xrightarrow{\sim}H^3(X,\Z).
\end{equation}

\begin{lemma}[Clean exceptional restriction]\label{lem:clean-base-change}
For \(T\in H^*(F\times F,\Z)\),
\begin{equation}\label{eq:clean-base-change}
 e^*q_*\mu^*T=p_*\pi^*\Delta^*T,
 \qquad
 b_*q_*\mu^*T=\psi_*T.
\end{equation}
\end{lemma}

\begin{proof}
The equality \(\psi^{-1}(0)=\Delta\) follows from injectivity of the
Albanese embedding.  Thus \(q^{-1}(X)=P\) set-theoretically.  It has
multiplicity one: on a fibre \(C\) of \(P\to F\), both \(q^*X\) and the
exceptional divisor have degree \(-1\).  Hence \(q^*X=P\) as Cartier
divisors and the normal bundles agree.  Naturality of Thom classes gives
\(q^*e_*=j_*p^*\), where \(j:P\hookrightarrow Y\).  The first identity of
\eqref{eq:clean-base-change} follows by adjunction and the projection
formula; the second follows from \(bq=\psi\mu\) and
\(\mu_*\mu^*=1\).
\end{proof}

We next calculate the endpoint of the difference correspondence.  The
Pontryagin product on \(J\) is denoted by \(\star\).

\begin{lemma}[Pontryagin endpoint]\label{lem:pontryagin-endpoint}
For a principally polarized abelian fivefold and
\(\xi\in\bigwedge^9\Lambda\), let \(y(\xi)\in\Lambda\) be its integral
symplectic wedge-dual.  With one global sign depending on orientation and
Pontryagin conventions,
\begin{equation}\label{eq:pontryagin-endpoint}
 \xi\star\theta^{[3]}=\pm\theta^{[2]}\wedge y(\xi).
\end{equation}
\end{lemma}

\begin{proof}
Choose a symplectic basis.  It suffices to take \(\xi\) to be the oriented
top monomial with one basis vector omitted.  In the adjunction formula
\[
 \int_J(\xi\star\theta^{[3]})\wedge\eta
 =\int_{J\times J}\operatorname{pr}_1^*\xi\wedge
 \operatorname{pr}_2^*\theta^{[3]}\wedge m^*\eta,
\]
the first factor must receive the omitted vector, while the second receives
the two complementary full symplectic pairs.  Each admissible term occurs
once because \(\theta^{[3]}\) is a divided power.  The result is the pairing
of \(\pm\theta^{[2]}\wedge y(\xi)\) with \(\eta\).  The degree-five wedge
pairing is unimodular, which proves the identity.
\end{proof}

For \(\beta\in H^3(F,\Z)/\operatorname{tors}\), put
\(\xi=a_*\beta\in H^9(J,\Z)\).  The Albanese isomorphism, Poincare duality,
and the principal polarization identify \(\beta\mapsto y(\xi)\) with an
integral isomorphism from the free part of \(H^3(F,\Z)\) to \(\Lambda\).
Combining it with the inverse of \eqref{eq:cylinder-adjoint}, and absorbing
the global sign, defines the isomorphism
\(y:H^3(X,\Z)\to\Lambda\) used in the introduction.

\begin{proposition}[Endpoint lifts]\label{prop:endpoint-lifts}
For every \(\gamma\in H^3(X,\Z)\), there is
\(u_\gamma\in H^3(M,\Z)\) such that
\begin{equation}\label{eq:endpoint-lift}
 e^*u_\gamma=\gamma,
 \qquad
 b_*u_\gamma=\theta^{[2]}\wedge y(\gamma).
\end{equation}
The common choice of sign is included in the definition of \(y\).
\end{proposition}

\begin{proof}
Choose \(\beta\) with \(p_*\pi^*\beta=\gamma\), and take
\(u_\gamma=q_*\mu^*(\beta\otimes1)\), exchanging the two factors if needed
to match the sign convention.  Lemma~\ref{lem:clean-base-change} gives the
restriction.  Since the difference map is addition composed with
\((-a)\times a\), the pushforward is
\[
 \psi_*(\beta\otimes1)=(-1)^*(a_*\beta)\star[F].
\]
Now \([F]=\theta^{[3]}\), and
Lemma~\ref{lem:pontryagin-endpoint} gives the second identity.
\end{proof}

Proposition~\ref{prop:endpoint-lifts} completes the proof of surjectivity in
Proposition~\ref{prop:topological-exact}.  Choose a basis of \(H^3(X,\Z)\)
and the corresponding lifts.  Together with
\(b^*\bigwedge^3\Lambda\), they form a basis of \(H^3(M,\Z)\).  Their
pushforwards are, respectively,
\(\theta^{[2]}\wedge\Lambda\) and
\[
 b_*b^*\alpha=\theta\wedge\alpha.
\]
Theorem~\ref{thm:symplectic-saturation} therefore proves
\eqref{eq:main-saturation}.

Finally, \((b_*,e^*)\) is injective.  Indeed, on \(\ker e^*=\im b^*\), its
first component is \(L\), and \(L\) is injective by the block calculation
in Section~\ref{sec:lattice}.  Changing a lift of \(\gamma\) adds
\(b^*\alpha\), hence changes its pushforward by \(L\alpha\).  Equation
\eqref{eq:endpoint-lift} and the isomorphism \eqref{eq:tau-isomorphism} now
give exactly the fibre product \eqref{eq:main-fibre-product}.  This proves
Theorem~\ref{thm:main-lattice}.
